\section{Conclusion}
\label{sec:conclusion}

Local green industrial policy is multidimensional. Governments can support a firm's own decarbonizing investment or improve shared green infrastructure, while the firms they host compete in broader product markets. This paper shows that the intensity of that competition can change not only the level but the direction of a cross-instrument policy response.

When firms endogenously choose both conventional cost-reducing investment and green investment, there are conditions under which a rival green subsidy reduces local infrastructure at low product substitutability but raises it once substitutability exceeds a unique threshold. The result follows from a strategic amplification channel: the rival subsidy improves the rival firm's competitiveness, business stealing changes the local firm's two investment choices, and the relative marginal value of targeted subsidy and shared infrastructure changes. A nested benchmark without conventional competitive investment does not reproduce the sign reversal.

The broader implication is that stronger interjurisdictional competition need not produce simple policy matching. It can instead induce governments to change the composition of their response toward instruments that jointly support environmental objectives and local competitiveness. This distinction is potentially useful for evaluating place-based green industrial policies and for designing empirical tests of strategic policy interaction across regions.
